\documentclass[a5paper,10pt,twoside]{article}

\usepackage[a4paper,inner=1.5cm,outer=1.5cm,top=2.0cm,bottom=2.0cm,pdftex]{geometry}

\usepackage[pdftex]{hyperref}
\hypersetup{
%%    bookmarks=false,         % show bookmarks bar?
    unicode=true,          % non-Latin characters in Acrobat’s bookmarks
%    pdftoolbar=true,        % show Acrobat’s toolbar?
%    pdfmenubar=true,        % show Acrobat’s menu?
%    pdffitwindow=true,      % page fit to window when opened
    pdftitle={Tiéd a dal},    % title
    pdfauthor={Kovács Tamás},     % author
%    pdfsubject={Gitáros kottás könyv},   % subject of the document
%    pdfnewwindow=true,      % links in new window
%    pdfkeywords={keywords}, % list of keywords
    colorlinks=true,       % false: boxed links; true: colored links
    linkcolor=black,          % color of internal links
    citecolor=black,        % color of links to bibliography
%    filecolor=magenta,      % color of file links
    urlcolor=black,           % color of external links
    pdflang=hu-HU,
    pdfpagelayout=TwoPageRight
}
% temp
%\newcommand{\url}[1]{{#1}}

\renewcommand{\ttdefault}{txtt}

\addto\captionsmagyar{\renewcommand{\indexname}{Tartalom}}

\newcommand{\cim}{Nótáskönyv}

\newcommand{\cimcentered}[2]
{
\begin{minipage}{\textwidth}
\vspace{0.4cm}
\begin{center}
\phantomsection
\addcontentsline{toc}{subsection}{#1}\label{#2}
\dalcim{#1}
\end{center}
\end{minipage}
}

\newcommand{\dalcim}[1]{\index{#1}\vspace{-0.1cm}\Large{\textbf{#1}}\vspace{0.0cm}}

\newcommand{\kottapart}[1]%
{
\begin{minipage}{\textwidth}
\vspace{0.0cm}%
\includegraphics[width=\textwidth]{generated/kotta/files/pdf/#1.pdf}%
\vspace{0.4cm}%
\end{minipage}
}

\newcommand{\kottastart}
{
\begin{minipage}{\textwidth}~
\vspace{0.5cm}
\end{minipage}
}

\newcommand{\kottaend}
{
\begin{minipage}{\textwidth}~
\vspace{0.3cm}
\end{minipage}
}


\newcommand{\noscorespacing}
{
\begin{minipage}{\textwidth}~
\vspace{0.3cm}
\end{minipage}
}

\newcommand{\versszakspacing}
{
\begin{minipage}{\textwidth}~
\vspace{0.3cm}
\end{minipage}
}

\newcommand{\alternatecim}[1]
{
\index{#1}
}

 

\begin{document}
\pagestyle{empty} 
\setcounter{page}{1}

\setcounter{tocdepth}{3}
\setcounter{secnumdepth}{2}       


\begin{titlepage}
%\setcounter{page}{1000}
\pagenumbering{roman}   % i, ii, iii, iv, ...
\begin{center}

    \begin{table}[ht]
        \begin{tabular*}{\textwidth}{@{\extracolsep{\fill}}lcr}
        \includegraphics[scale=0.14]{pictures/liliomv_black.pdf} &  & \includegraphics[scale=0.28]{pictures/ivbela_black.pdf}\\
        \end{tabular*}
    \end{table}%





{\large \cscslong}
%\vspace{0.1cm}
%~\\[\baselineskip]
%~\\[\baselineskip]
%~\\[\baselineskip]


\includegraphics[scale=0.45, trim = 15mm 15mm 5mm 15mm, clip]{pictures/magyarnepmesek.pdf}
%~\\[\baselineskip]
%~\\[\baselineskip]


%~\\[\baselineskip]

\vfill

\end{center}

\begin{center}


{\large \textsc{\kiado}}

\end{center}
 
\end{titlepage}
%\pagenumbering{alph}


\newpage
~\newpage 

\begin{center}

\mbox{}
\newpage

    \begin{table}[ht]
        \begin{tabular*}{\textwidth}{@{\extracolsep{\fill}}lcr}
        \includegraphics[scale=0.197]{pictures/liliomv_black.pdf} &  & \includegraphics[scale=0.39]{pictures/ivbela_black.pdf}\\
%        \includegraphics[width=0.2\textwidth]{pictures/liliomv_black.pdf} &  & \includegraphics[width=0.2\textwidth]{pictures/ivbela_black.pdf}\\
        \end{tabular*}
    \end{table}%

~\\[\baselineskip]
~\\[\baselineskip]

{\Huge \textbf{\cim}}
~\\[\baselineskip]
~\\[\baselineskip]
{\large \cscslong}
~\\[\baselineskip]
~\\[\baselineskip]
%Laura példánya. basszuskulcsban :)
\vfill

%{\large \textsc{\kiado}}
\end{center}

\newpage
%\mbox{}
%\newpage

\begin{center}
{\large
Szerkesztette:

Kovács Tamás

\begin{minipage}{\textwidth}
\vspace{0.6cm}
\end{minipage}

A borítót és a csapatlogót tervezte:

Bakos András

\begin{minipage}{\textwidth}
\vspace{0.6cm}
\end{minipage}

Zeneileg  szerkesztette:

Vörös Rózsa}
%\end{center}

~\\[\baselineskip]
~\\[\baselineskip]
~\\[\baselineskip]
~\\[\baselineskip]
~\\[\baselineskip]
~\\[\baselineskip]
{\large

Összesen 178dal található a könyvben.

A könyv \LaTeX~és Lilypond segítségével készült, Mediawikiből php-vel fordítva.

Az észrevételeket a \mbox{\href{mailto:kovianyo@gmail.com}{kovianyo@gmail.com}} címre várjuk.

Köszönet Regéné Viktor Zsuzsinak a második kiadás előtti alapos átnézésért.

A kézirat lezárva: \today

\begin{minipage}{\textwidth}
\vspace{0.6cm}
\end{minipage}

\url{http://105.cserkesz.hu/}
}
\vfill

%\begin{center}
{\large
Második kiadás

\the\year}
\end{center}


\newpage

\begin{center}

~\\[\baselineskip]
~\\[\baselineskip]%

%{\LARGE \cscs}
~\\[\baselineskip] %

{\Huge \textbf{\cim}} 
~\\[\baselineskip]
~\\[\baselineskip]
{\large A gitáros zenekar dalai}

\end{center}

~\\[\baselineskip] %
~\\[\baselineskip]
~\\[\baselineskip]
\begin{center}
\begin{minipage}{0.5\textwidth}
{
%,,
\begin{alltt}\normalfont\large{}
Megmondom a titkát, édesem a dalnak:
Önmagát hallgatja, aki dalra hallgat.
Mindenik embernek a lelkében dal van,
és a saját lelkét hallja minden dalban.
És akinek szép a lelkében az ének,
az hallja a mások énekét is szépnek.

                \textit{Babits Mihály: A második ének}
\end{alltt}
%''
~\\[\baselineskip] %
~\\[\baselineskip] %
~\\[\baselineskip]
~\\[\baselineskip]
}
\end{minipage}
\end{center}
{\large
Zenekarunk két fővel alakult meg 1990-ben. Az évek során sok fiatal csatlakozott az együtteshez, és mára közel húszan zenélünk együtt vasárnaponként a gitáros misén. Az elmúlt húsz évben rengeteg dalt énekeltünk, amelyeket ebbe a kottafüzetbe gyűjtöttünk össze. 

Szeretnénk ajánlani ezt a gyűjteményt azoknak, akik szívesen zenélnek és énekelnek velünk. Mindenkinek kívánjuk, hogy fedezze fel benne kedvenceit, hisz \cim.

}

\begin{minipage}{\textwidth}
\vspace{0.8cm}
\raggedleft{\large{\textit{Rózsa}}}
\end{minipage}
\normalsize

\newpage
\mbox{}
\newpage




\large

\pagestyle{fancy}
%\renewcommand{\sectionmark}[1]{\markboth{Dalsokönyv}{}}

%\renewcommand{\subsectionmark}[1]{\markright{#1}}
%\renewcommand{\subsectionmark}[1]{\markright{\thesubsection.\ #1}{}}
%http://texblog.wordpress.com/2007/11/07/headerfooter-in-latex-with-fancyhdr/
%\setlength\headheight{19pt}
\resetheaders

\newpage


\pagenumbering{arabic}  % 1, 2, 3, 4, ...
\setcounter{page}{1} % oldal számláló 1-re állítása

% from now on it is twosided
% http://www.latex-community.org/forum/viewtopic.php?f=4&t=38
\setboolean{@twoside}{true} 
% maxSorLength: 0
% virtualLines: 10
% kottaSorCount: 2
% count(lines): 1
% columnCount: 1
\begin{minipage}{\textwidth}
\cimcentered{Ablakomba besütött a holdvilág}{AblakombabesC3BCtC3B6ttaholdvilC3A1g}
\kottastart
\kottapart{AblakombabesC3BCtC3B6ttaholdvilC3A1g-1}
\kottapart{AblakombabesC3BCtC3B6ttaholdvilC3A1g-2}
\kottaend
\end{minipage}
~\vspace{1.0cm}
\newline
%
\begin{minipage}{\textwidth}
\cimcentered{A cigányok sátora}{AcigC3A1nyoksC3A1tora}
\kottastart
\kottapart{AcigC3A1nyoksC3A1tora-1}
\kottapart{AcigC3A1nyoksC3A1tora-2}
\kottapart{AcigC3A1nyoksC3A1tora-3}
\kottaend
\end{minipage}
~\vspace{1.0cm}
\newline
%
% maxSorLength: 51
% virtualLines: 10
% kottaSorCount: 2
% count(lines): 14
% columnCount: 1
\begin{minipage}{\textwidth}
\cimcentered{A csitári hegyek alatt}{AcsitC3A1rihegyekalatt}
\kottastart
\kottapart{AcsitC3A1rihegyekalatt-1}
\kottapart{AcsitC3A1rihegyekalatt-2}
\kottaend
\end{minipage}
~\vspace{1.0cm}
\newline
%


% remove thumb indexes from header

%\pagestyle{empty}
%\begin{landscape}
%\include{generated/tablazat}
%\end{landscape}  
%\index{szavak}{szó}
\newpage

\resetheaders

\setlength\parindent{0.3in}
\normalsize
\section*{Bevezető}
\input{generated/files/bevezeto}

\newpage
\setlength\parindent{0in}
\begin{multicols}{2} 
\section*{Dalok tematikusan}
\footnotesize Szent, szent a neve \ \textbf{\pageref{Szent2Cszentaneve}} \newline
\footnotesize Napfivér, Holdnővér \ \textbf{\pageref{NapfivC3A9r2CHoldnC591vC3A9r}} \newline
\footnotesize Minden, mi él \ \textbf{\pageref{Minden2CmiC3A9l}} \newline
\footnotesize Kenyered és borod \ \textbf{\pageref{KenyeredC3A9sborod}} \newline
\footnotesize Jöjj hozzám \ \textbf{\pageref{JC3B6jjhozzC3A1m}} \newline
\footnotesize Megszentelt kenyér \ \textbf{\pageref{MegszenteltkenyC3A9r}} \newline
\footnotesize Ő az Út, az Élet \ \textbf{\pageref{C590azC39At2CazC389let}} \newline
\footnotesize Lelkem, az Urat dicsérd \ \textbf{\pageref{Lelkem2CazUratdicsC3A9rd}} \newline
\footnotesize Ó, jöjjetek, örvendjünk \ \textbf{\pageref{C3932CjC3B6jjetek2CC3B6rvendjC3BCnk}} \newline
\footnotesize Hajlékomba \ \textbf{\pageref{HajlC3A9komba}} \newline
\footnotesize Békesség legyen velünk \ \textbf{\pageref{BC3A9kessC3A9glegyenvelC3BCnk}} \newline
\footnotesize A szeretet mindenkié \ \textbf{\pageref{AszeretetmindenkiC3A9}} \newline
\footnotesize Isten Báránya új \ \textbf{\pageref{IstenBC3A1rC3A1nyaC3BAj}} \newline
\footnotesize Gyönyörű vagy, én Uram \ \textbf{\pageref{GyC3B6nyC3B6rC5B1vagy2CC3A9nUram}} \newline
\footnotesize Atyám, Tied vagyok \ \textbf{\pageref{AtyC3A1m2CTiedvagyok}} \newline
\footnotesize Királyi törvény \ \textbf{\pageref{KirC3A1lyitC3B6rvC3A9ny}} \newline
\footnotesize Irgalmazz, Uram, irgalmazz \ \textbf{\pageref{Irgalmazz2CUram2Cirgalmazz}} \newline
\footnotesize Kegyelmed, Úr Jézus \ \textbf{\pageref{Kegyelmed2CC39ArJC3A9zus}} \newline
\footnotesize Fölém hajló jóságodra \ \textbf{\pageref{FC3B6lC3A9mhajlC3B3jC3B3sC3A1godra}} \newline
\footnotesize Szent, szent, szent az Isten \ \textbf{\pageref{Szent2Cszent2CszentazIsten}} \newline
\footnotesize Glória \ \textbf{\pageref{GlC3B3ria}} \newline
\footnotesize Köszönöm, Jézus \ \textbf{\pageref{KC3B6szC3B6nC3B6m2CJC3A9zus}} \newline
\footnotesize Uram, irgalmazz \ \textbf{\pageref{Uram2Cirgalmazz}} \newline
\footnotesize Csak állok szótlanul \ \textbf{\pageref{CsakC3A1llokszC3B3tlanul}} \newline
\footnotesize Király vagy \ \textbf{\pageref{KirC3A1lyvagy}} \newline
\footnotesize Krisztus az Úr \ \textbf{\pageref{KrisztusazC39Ar}} \newline
\footnotesize Menjetek be kapuin \ \textbf{\pageref{Menjetekbekapuin}} \newline
\footnotesize Jöjj, Szentlélek \ \textbf{\pageref{JC3B6jj2CSzentlC3A9lek}} \newline
\footnotesize Isten Báránya régi \ \textbf{\pageref{IstenBC3A1rC3A1nyarC3A9gi}} \newline
\footnotesize Te vagy az út, az igazság \ \textbf{\pageref{TevagyazC3BAt2CazigazsC3A1g}} \newline
\footnotesize Teremts bennem \ \textbf{\pageref{Teremtsbennem}} \newline
\footnotesize Jézus, Te közénk eljöttél \ \textbf{\pageref{JC3A9zus2CTekC3B6zC3A9nkeljC3B6ttC3A9l}} \newline
\footnotesize Nagy Úr \ \textbf{\pageref{NagyC39Ar}} \newline
\footnotesize Rád bízom \ \textbf{\pageref{RC3A1dbC3ADzom}} \newline
\footnotesize Mint szarvas hűs \ \textbf{\pageref{MintszarvashC5B1s}} \newline
\footnotesize Isten szeret bennünket \ \textbf{\pageref{IstenszeretbennC3BCnket}} \newline
\footnotesize Nem lesz egyedül a szívem \ \textbf{\pageref{NemleszegyedC3BClaszC3ADvem}} \newline
\footnotesize Tüzed, Uram, Jézus \ \textbf{\pageref{TC3BCzed2CUram2CJC3A9zus}} \newline
\footnotesize Csak Benned \ \textbf{\pageref{CsakBenned}} \newline
\footnotesize Ároni áldás \ \textbf{\pageref{C381roniC3A1ldC3A1s}} \newline
\footnotesize Ez az a nap \ \textbf{\pageref{Ezazanap}} \newline
\footnotesize Együtt fog ujjongni \ \textbf{\pageref{EgyC3BCttfogujjongni}} \newline
\footnotesize Róla beszél fű, virág \ \textbf{\pageref{RC3B3labeszC3A9lfC5B12CvirC3A1g}} \newline
\footnotesize Jézussal járok mindennap \ \textbf{\pageref{JC3A9zussaljC3A1rokmindennap}} \newline
\footnotesize Áldalak, jó Uram \ \textbf{\pageref{C381ldalak2CjC3B3Uram}} \newline
\footnotesize Áldásoddal megyünk \ \textbf{\pageref{C381ldC3A1soddalmegyC3BCnk}} \newline
\footnotesize Ízleld és lásd, hogy jó az Úr \ \textbf{\pageref{C38DzleldC3A9slC3A1sd2ChogyjC3B3azC39Ar}} \newline
\footnotesize Szentlélek \ \textbf{\pageref{SzentlC3A9lek}} \newline
\footnotesize Kész az én szívem \ \textbf{\pageref{KC3A9szazC3A9nszC3ADvem}} \newline
\footnotesize Szállj, dalom \ \textbf{\pageref{SzC3A1llj2Cdalom}} \newline
\footnotesize Tied a dicsőség \ \textbf{\pageref{TiedadicsC591sC3A9g}} \newline
\footnotesize Isten a szépség, Isten a jó \ \textbf{\pageref{IstenaszC3A9psC3A9g2CIstenajC3B3}} \newline
\footnotesize Röpke fohász \ \textbf{\pageref{RC3B6pkefohC3A1sz}} \newline
\footnotesize Itt vagyok végre \ \textbf{\pageref{IttvagyokvC3A9gre}} \newline
\footnotesize Mondjunk éneket örömmel \ \textbf{\pageref{MondjunkC3A9neketC3B6rC3B6mmel}} \newline
\footnotesize Szeretem örökké \ \textbf{\pageref{SzeretemC3B6rC3B6kkC3A9}} \newline
\footnotesize Csak Jézusnak szolgálok \ \textbf{\pageref{CsakJC3A9zusnakszolgC3A1lok}} \newline
\footnotesize Dicsérjétek az Urat \ \textbf{\pageref{DicsC3A9rjC3A9tekazUrat}} \newline
\footnotesize Életem csak kegyelem \ \textbf{\pageref{C389letemcsakkegyelem}} \newline
\footnotesize Áldott Istenünk \ \textbf{\pageref{C381ldottIstenC3BCnk}} \newline
\footnotesize Hálát adok, hogy itt a reggel \ \textbf{\pageref{HC3A1lC3A1tadok2Chogyittareggel}} \newline
\footnotesize Hálát adok az esti órán \ \textbf{\pageref{HC3A1lC3A1tadokazestiC3B3rC3A1n}} \newline
\footnotesize Bárányom \ \textbf{\pageref{BC3A1rC3A1nyom}} \newline
\footnotesize Áldjad, én lelkem, az Urat \ \textbf{\pageref{C381ldjad2CC3A9nlelkem2CazUrat}} \newline
\footnotesize Jézus, nékünk adtál anyát \ \textbf{\pageref{JC3A9zus2CnC3A9kC3BCnkadtC3A1lanyC3A1t}} \newline
\footnotesize Jól vigyázz, kicsi kéz \ \textbf{\pageref{JC3B3lvigyC3A1zz2CkicsikC3A9z}} \newline
\footnotesize Mária, jó anyánk \ \textbf{\pageref{MC3A1ria2CjC3B3anyC3A1nk}} \newline
\footnotesize Mária, Mária, szép virágszál \ \textbf{\pageref{MC3A1ria2CMC3A1ria2CszC3A9pvirC3A1gszC3A1l}} \newline
\footnotesize Ő az Úr, Ő az Úr \ \textbf{\pageref{C590azC39Ar2CC590azC39Ar}} \newline
\footnotesize Legyen a béke teveled \ \textbf{\pageref{LegyenabC3A9keteveled}} \newline
\footnotesize Nap, hold csillagok \ \textbf{\pageref{Nap2Choldcsillagok}} \newline
\footnotesize Napkeltétől \ \textbf{\pageref{NapkeltC3A9tC591l}} \newline
\footnotesize Szeretet áradjon köztünk \ \textbf{\pageref{SzeretetC3A1radjonkC3B6ztC3BCnk}} \newline
\footnotesize Szent vagy, mindenség Ura, Istene \ \textbf{\pageref{Szentvagy2CmindensC3A9gUra2CIstene}} \newline
\footnotesize Szent vagy, szent vagy, mindenség Ura \ \textbf{\pageref{Szentvagy2Cszentvagy2CmindensC3A9gUra}} \newline
\footnotesize Táncolj az Úrnak, dicsérjed nevét \ \textbf{\pageref{TC3A1ncoljazC39Arnak2CdicsC3A9rjednevC3A9t}} \newline
\footnotesize Add a kezed \ \textbf{\pageref{Addakezed}} \newline
\footnotesize Terád vár egy szép ország \ \textbf{\pageref{TerC3A1dvC3A1regyszC3A9porszC3A1g}} \newline
\footnotesize Zengd velünk \ \textbf{\pageref{ZengdvelC3BCnk}} \newline
\footnotesize Ne félj, ne aggódj \ \textbf{\pageref{NefC3A9lj2CneaggC3B3dj}} \newline
\footnotesize Irgalmas Istenünk jóságát \ \textbf{\pageref{IrgalmasIstenC3BCnkjC3B3sC3A1gC3A1t}} \newline
\footnotesize Várj és ne félj \ \textbf{\pageref{VC3A1rjC3A9snefC3A9lj}} \newline
\footnotesize Ahol szeretet és jóság \ \textbf{\pageref{AholszeretetC3A9sjC3B3sC3A1g}} \newline
\footnotesize Ó, halld meg, Uram \ \textbf{\pageref{C3932Challdmeg2CUram}} \newline
\footnotesize Jó az Úrban bizakodni \ \textbf{\pageref{JC3B3azC39Arbanbizakodni}} \newline
\footnotesize Jézus életem, erőm, békém \ \textbf{\pageref{JC3A9zusC3A9letem2CerC591m2CbC3A9kC3A9m}} \newline
\footnotesize Bizakodjatok, jó az Úr \ \textbf{\pageref{Bizakodjatok2CjC3B3azC39Ar}} \newline
\footnotesize Áldott légy, Uram \ \textbf{\pageref{C381ldottlC3A9gy2CUram}} \newline
\footnotesize Isteni béke \ \textbf{\pageref{IstenibC3A9ke}} \newline
\footnotesize A béke napja \ \textbf{\pageref{AbC3A9kenapja}} \newline
\footnotesize A föld és az ég \ \textbf{\pageref{AfC3B6ldC3A9sazC3A9g}} \newline
\footnotesize Ím néped \ \textbf{\pageref{C38DmnC3A9ped}} \newline
\footnotesize A mélyből Hozzád száll szavam \ \textbf{\pageref{AmC3A9lybC591lHozzC3A1dszC3A1llszavam}} \newline
\footnotesize Föltámadt Jézus \ \textbf{\pageref{FC3B6ltC3A1madtJC3A9zus}} \newline
\footnotesize Érzem, hogy hinnem kell \ \textbf{\pageref{C389rzem2Chogyhinnemkell}} \newline
\footnotesize Nyújtsd a kezed \ \textbf{\pageref{NyC3BAjtsdakezed}} \newline
\footnotesize Múlnak az évek \ \textbf{\pageref{MC3BAlnakazC3A9vek}} \newline
\footnotesize Ne félj, mert megváltottalak \ \textbf{\pageref{NefC3A9lj2CmertmegvC3A1ltottalak}} \newline
\footnotesize Tápláló, szent kenyér \ \textbf{\pageref{TC3A1plC3A1lC3B32CszentkenyC3A9r}} \newline
\footnotesize Látod, újra este van \ \textbf{\pageref{LC3A1tod2CC3BAjraestevan}} \newline
\footnotesize Ha jön az Úr \ \textbf{\pageref{HajC3B6nazC39Ar}} \newline
\footnotesize Menj, Ábrahám \ \textbf{\pageref{Menj2CC381brahC3A1m}} \newline
\footnotesize Nézz, testvér, fel \ \textbf{\pageref{NC3A9zz2CtestvC3A9r2Cfel}} \newline
\footnotesize Szeretlek, Uram \ \textbf{\pageref{Szeretlek2CUram}} \newline
\footnotesize Álmaimból, lelkemből \ \textbf{\pageref{C381lmaimbC3B3l2ClelkembC591l}} \newline
\footnotesize Itt a két kezem \ \textbf{\pageref{IttakC3A9tkezem}} \newline
\footnotesize Gyertyaláng \ \textbf{\pageref{GyertyalC3A1ng}} \newline
\footnotesize Pásztoroknak pásztora \ \textbf{\pageref{PC3A1sztoroknakpC3A1sztora}} \newline
\footnotesize Krisztus Urunknak áldott születésén \ \textbf{\pageref{KrisztusUrunknakC3A1ldottszC3BCletC3A9sC3A9n}} \newline
\footnotesize Mikulás \ \textbf{\pageref{MikulC3A1s}} \newline
\footnotesize Betlehem, Betlehem \ \textbf{\pageref{Betlehem2CBetlehem}} \newline
\footnotesize A kis Jézus megszületett \ \textbf{\pageref{AkisJC3A9zusmegszC3BCletett}} \newline
\footnotesize Ó, gyönyörű szép \ \textbf{\pageref{C3932CgyC3B6nyC3B6rC5B1szC3A9p}} \newline
\footnotesize Mennyből az angyal \ \textbf{\pageref{MennybC591lazangyal}} \newline
\footnotesize Kit sok boldog pásztor \ \textbf{\pageref{KitsokboldogpC3A1sztor}} \newline
\footnotesize Hadd zengjen énekszó \ \textbf{\pageref{HaddzengjenC3A9nekszC3B3}} \newline
\footnotesize Karácsony \ \textbf{\pageref{KarC3A1csony}} \newline
\footnotesize Az angyal énekel \ \textbf{\pageref{AzangyalC3A9nekel}} \newline
\footnotesize Fel nagy örömre \ \textbf{\pageref{FelnagyC3B6rC3B6mre}} \newline
\footnotesize Advent \ \textbf{\pageref{Advent}} \newline
\footnotesize Pásztorok, pásztorok \ \textbf{\pageref{PC3A1sztorok2CpC3A1sztorok}} \newline
\footnotesize Az ősi jóslat \ \textbf{\pageref{AzC591sijC3B3slat}} \newline
\footnotesize Kirje, kirje \ \textbf{\pageref{Kirje2Ckirje}} \newline
\footnotesize Csendes éj \ \textbf{\pageref{CsendesC3A9j}} \newline
\footnotesize Hóban, fényben \ \textbf{\pageref{HC3B3ban2CfC3A9nyben}} \newline
\footnotesize Évezredek óta \ \textbf{\pageref{C389vezredekC3B3ta}} \newline
\footnotesize Mária altatója \ \textbf{\pageref{MC3A1riaaltatC3B3ja}} \newline
\footnotesize Szálljatok le \ \textbf{\pageref{SzC3A1lljatokle}} \newline
\footnotesize Jöjj el, Betlehem \ \textbf{\pageref{JC3B6jjel2CBetlehem}} \newline
\footnotesize Karácsony ünnep \ \textbf{\pageref{KarC3A1csonyC3BCnnep}} \newline
\footnotesize Menedékem az Isten \ \textbf{\pageref{MenedC3A9kemazIsten}} \newline
\footnotesize Jézus, én Téged választottalak \ \textbf{\pageref{JC3A9zus2CC3A9nTC3A9gedvC3A1lasztottalak}} \newline
\footnotesize Én Uram, én Istenem \ \textbf{\pageref{C389nUram2CC3A9nIstenem}} \newline
\footnotesize Álmodjék Rólad a szívünk \ \textbf{\pageref{C381lmodjC3A9kRC3B3ladaszC3ADvC3BCnk}} \newline
\footnotesize Ó, jöjj, ó, jöjj, Üdvözítő \ \textbf{\pageref{C3932CjC3B6jj2CC3B32CjC3B6jj2CC39CdvC3B6zC3ADtC591}} \newline
\footnotesize Istentől elküldetett \ \textbf{\pageref{IstentC591lelkC3BCldetett}} \newline
\footnotesize Millenniumi óda \ \textbf{\pageref{MillenniumiC3B3da}} \newline
\footnotesize Rossz a Jézus kiscsizmája \ \textbf{\pageref{RosszaJC3A9zuskiscsizmC3A1ja}} \newline
\footnotesize Betlehem városban \ \textbf{\pageref{BetlehemvC3A1rosban}} \newline
\footnotesize Angyal szárnya \ \textbf{\pageref{AngyalszC3A1rnya}} \newline
\footnotesize Az volna jó \ \textbf{\pageref{AzvolnajC3B3}} \newline
\footnotesize Jézus, köszönöm Néked \ \textbf{\pageref{JC3A9zus2CkC3B6szC3B6nC3B6mNC3A9ked}} \newline
\footnotesize Ég az a tűz \ \textbf{\pageref{C389gazatC5B1z}} \newline
\footnotesize Neked hála \ \textbf{\pageref{NekedhC3A1la}} \newline
\footnotesize Jézus ágyán nincsen paplan \ \textbf{\pageref{JC3A9zusC3A1gyC3A1nnincsenpaplan}} \newline
\footnotesize Betlehem kis falucskába \ \textbf{\pageref{BetlehemkisfalucskC3A1ba}} \newline
\footnotesize A nap bíborban áll \ \textbf{\pageref{AnapbC3ADborbanC3A1ll}} \newline
\footnotesize Esti köszöntő \ \textbf{\pageref{EstikC3B6szC3B6ntC591}} \newline
\footnotesize Karácsonyi álom \ \textbf{\pageref{KarC3A1csonyiC3A1lom}} \newline
\footnotesize Esik a hó \ \textbf{\pageref{EsikahC3B3}} \newline
\footnotesize Misztrál Advent \ \textbf{\pageref{MisztrC3A1lAdvent}} \newline
\footnotesize Kicsi vagyok én \ \textbf{\pageref{KicsivagyokC3A9n}} \newline
\footnotesize Életünk kaptuk Tőled, Uram \ \textbf{\pageref{C389letC3BCnkkaptukTC591led2CUram}} \newline

\end{multicols}
\newpage


% kották jegyzéke
%\listoffigures
%\newpage

\fancyfoot[er,ol]{}
\fancyhead[el,or]{\includegraphics[scale=0.05]{pictures/ivbela_black.pdf}}

\footnotesize
\begin{multicols}{2} 
\tableofcontents
\end{multicols}

%\newpage

%\fancyhead[r]{Tárgymutató}
%\footnotesize 
%\printindex

%\printindex{szavak}{Szószedet}

%\pagestyle{empty} 


\pagestyle{empty} 
%\newpage
%~
\newpage

\fancyfoot[lr]{}
\fancyhead[lr]{}
\renewcommand{\headrulewidth}{0pt}

\pagestyle{fancy}
\setcounter{betu}{0}
\addtocounter{betu}{1}\extramarks{A}{}
\fancyhead[el]{\overview%\lblob\addtocounter{betu}{1}\extramarks{B}{}\lblob
}

~
\newpage
~
\newpage

%\fancyhead[or]{\includegraphics[scale=0.05]{pictures/ivbela_black.pdf}\rblob}

\pagestyle{empty} 

\newpage

% trim=l b r t
~\newline
\begin{center}
\includegraphics[scale=0.6, angle=0, trim = 10mm 20mm 10mm 0mm]{pictures/csapatpolo_hatulja_black.pdf}
\end{center}

\end{document}



